\section{Holonomic conditions of the general gravitational selector}

The second gravitational template starts from
\begin{equation}
\label{rm:eq:gravity-delta-theta}
\delta_\theta
=\frac{A_{\rm deg}-C^*_{\rm deg}}{360}
\end{equation}
and a route length
\(\Lambda_{27}=\mathcal L_{\rm hol}(\widetilde{\mathcal K}_{27})\)
in an enriched lift of the return. The expression
\begin{equation}
\label{rm:eq:gravity-holonomic-template}
\boxed{
G_{\rm hol}
=\frac{c_{\rm int}^3}{\hbar_{\rm int}}
\bigl(\delta_\theta\Lambda_{27}\bigr)^2}
\end{equation}
is dimensionally exact. Its promotion to a general selector requires that
\(\mathcal L_{\rm hol}\) be natural under cyclic change of basis,
conjugation, relabeling, and subdivision, and possess a normalization
that is unique and stable under refinement.

There is additionally a torsion control: every positive multiplicative
character \(\chi:K_{27}\to\R_+^\times\) on a finite observable cycle of order
four is trivial, because \(\chi(g)^4=1\) and \(\chi(g)>0\) imply
\(\chi(g)=1\). A nontrivial holonomic response must therefore
be defined on the nontorsional lift or on enriched
memory; it cannot
be obtained by renaming a positive character of the finite cycle.

The chart
\begin{equation}
\label{rm:eq:gravity-decimal-chart}
10^{-11}\left(6+\frac{674}{1000}+\frac{30}{10^5}\right)
=6.67430\times10^{-11}
\end{equation}
is an exact evaluation of declared inputs. It does not fix
\(\Lambda_{27}\) or prove that \cref{rm:eq:gravity-holonomic-template}
assumes that value.

\begin{keyresult}
The periodicity condition determines \(G\) without using its decimal
expansion: the chronological structure CT108 selects \((54/\pi)^2\), the
gravitational line fixes the dimensional type, and the bidirectional family
preserves \(G\hbar\). The resulting coupling is represented on the
tensorial carrier constructed by \(P_5\), whereas the reduction to the
physical trace-free transverse sector is effected by a distinct
rank-two map.
\end{keyresult}

\begin{falsifier}
The circular block is refuted if
\(r_g(\theta)\neq\theta(\pi/54)\ell_0\), if there exists another positive solution
of \cref{rm:eq:gravity-closing-condition}, or if
\(G_a\hbar_a/c^3\neq\theta_{108}^{\rm circ}\ell_0^2\). The tensorial block is
refuted if \(\pi_{\rm ico}\) does not preserve incidence or intertwine the
declared action, if its stabilizer is not \(C_5\), if \(P_5\) is not a
self-adjoint idempotent of trace five, if
\(\Gamma_0\Gamma_0^*\neq(8/5)I\), or if \(\Lambda(n)\) does not have rank two.
The five-polarization reading is refuted if the tensors of
\cref{rm:eq:gravity-helicity2-real,rm:eq:gravity-helicity1-real,rm:eq:gravity-helicity0-real}
are not an orthonormal basis or if \(\Lambda(n)\) retains a weight other than
\(\pm2\).
The general holonomic selector is resolved by obstruction: without a
natural, normalized, and stable \(\mathcal L_{\rm hol}\) there is no
univocal selection compatible with all relabelings; using
\cref{rm:eq:gravity-decimal-chart} to choose it incurs circularity and
invalidates the construction.
\end{falsifier}

\begin{statusbox}
Transducer, uniqueness of the selector under the condition of
CT108 periodicity, and Planck family: \status{Exact internal theorem under the stated geometric condition}.
The descent \(\Omega_{12}^{\rm ico}\simeq A_5/C_5\) is exact under the four hypotheses stated explicitly in \(\pi_{\rm ico}\); the HMT construction of that map belongs to the exceptional provenance corridor. Once the descent, projector \(P_5\), intertwiner \(\Gamma_5\), five-helicity basis, TT reduction, and spectrum of \(\mathbb G_5\) are fixed: \status{Internal exact theorem}.
Massive reading, field equations, and realization
of Palatini--Cartan--Holst: \status{Realization, Physically Typed}. Template
holonomic: \status{Proved in homogeneity}; the nonexistence of a general unambiguous selector without additional normalization is a
\status{obstruction of canonicity proved}. Carta SI:
\status{Contrast metrological external}.
\end{statusbox}
